\documentclass[11pt]{article}
\usepackage[margin=1in]{geometry}
\usepackage{amsmath}
\usepackage{booktabs}
\usepackage[hidelinks]{hyperref}
\emergencystretch=3em

\newcommand{\code}[1]{\texttt{#1}}

\title{What the AF10R1 user guide changes in the accelerometer timing analysis}
\author{Ilan Upfal}
\date{28 October 2026}

\begin{document}
\maketitle

\section*{Summary}

This note follows up the report \emph{Accelerometer sample timing: dropped
samples and biased frequencies} (\code{reports/accel\_sampling}) using the
vendor's user guide for the accelerometer board (InUnit AB, AF10R1 User Guide,
DOC-2026-06-12-01 rev v1, draft; \code{data/pcb\_documentation/AF10R1\_UG.pdf}).
It changes three things.

\begin{enumerate}
  \item \textbf{The sensors sample at about 1493\,Hz, not 1600\,Hz.} We lose
        about 16\% of samples, not 22\%. The 62.7\,Hz first flap mode found
        by the corrected method is unaffected. The corrected higher-mode
        peaks in that report are not: they move by about 15\,Hz, and the
        second flap mode now matches the video's 164.5\,Hz.
  \item \textbf{The board can keep up.} The vendor streamed for 5 minutes with
        the same 84-set reads and lost no samples. Our read loop is the
        bottleneck.
  \item \textbf{Dropped samples are not buffered.} The guide describes a
        1024-set buffer, which should hold the first $\sim$4\,s of
        shortfall. Our data show samples missing from the start. The most
        likely cause is that a read discards the samples it does not return.
        Logging three registers would confirm it.
\end{enumerate}

\section{The actual sample rate is about 1493\,Hz}

The LIS2DS12 sensors run on an internal RC oscillator. The guide measured
1493\,Hz over 5 minutes (stable to $\pm$1\,Hz, all four sensors), 6.7\%
below the configured 1600\,Hz. Its specification table says 1525\,Hz. The
previous report assumed exactly 1600\,Hz.

\paragraph{Our data agree.} With the corrected timing (method~3 of the
previous report: samples on a fixed grid within each batch, gaps masked),
copies of the first mode appear 14.9\,Hz on either side of it. If the
assumed rate is too high, the copy above is stronger than the copy below.
Table~\ref{tab:copies} shows this lopsidedness at 1600\,Hz. The two copies are
about equal between 1493 and 1525\,Hz, the range the guide gives.

\begin{table}[h]
  \centering
  \begin{tabular}{lccccc}
    \toprule
    Assumed rate (Hz) & 1475 & 1493 & 1510 & 1525 & 1600 \\
    \midrule
    Copy at $f_1 - 14.9$\,Hz & 0.071 & 0.044 & 0.034 & 0.024 & 0.007 \\
    Copy at $f_1 + 14.9$\,Hz & 0.014 & 0.023 & 0.032 & 0.039 & 0.185 \\
    First mode $f_1$ (Hz) & 62.5 & 62.5 & 62.5 & 62.5 & 62.5 \\
    \bottomrule
  \end{tabular}
  \caption{Power of the batch-gap copies relative to the first flap mode
  (median PSD of \code{s4\_z}, 30 flap trials) for different assumed
  sensor rates.}
  \label{tab:copies}
\end{table}

\paragraph{Revised loss numbers.} An 84-sample batch takes 56.3\,ms to
produce, not 52.5\,ms. In each 67.3\,ms read cycle the sensor produces
about 100.5 samples, and we keep 84. The delivered rate of 1249\,Hz is 84\%
of the true rate, so we lose 16\%. The conclusion that samples are dropped
between batches still holds.

\paragraph{What this means for frequencies.} The batch read times fix which
frequencies the spectrum can show: the true frequency and its copies
$\pm n \times 14.9$\,Hz. The assumed rate only decides which of these
appears strongest. For the first mode the true line wins at every rate from
1475 to 1600\,Hz, so 62.5\,Hz stands. The two wrong-timing predictions in the
previous report (47.9 and 61.4\,Hz) also come out the same at 1493\,Hz.

For higher modes, a 7\% error in the rate moves the strongest line by one or
more copies. Table~\ref{tab:peaks} repeats the higher peaks at the corrected
rate.

\begin{table}[h]
  \centering
  \begin{tabular}{llcc}
    \toprule
    Excitation & Channel & Previous report (1600\,Hz) & Assumed 1493--1525\,Hz \\
    \midrule
    Flap & \code{s4\_z} & 176--178 & 159.5--162 \quad(video: 164.5) \\
    Flap & \code{s4\_z} & 275--285 & 260--267 \\
    Flap & \code{s4\_z} & $\approx$430 & 407--423 \\
    Edge & \code{s4\_y} & 118--121 & 115.5--116.5 \\
    Edge & \code{s4\_y} & 313--326 & 297--301 \\
    Edge & \code{s4\_y} & 577--597 & 534--559 \\
    \bottomrule
  \end{tabular}
  \caption{Higher peaks of the median PSDs. Right column: range of peak
  locations as the assumed rate varies from 1493 to 1525\,Hz. Previous
  values are over several sensors; new values are for sensor~4 only.}
  \label{tab:peaks}
\end{table}

The spread in the right-hand column is the uncertainty: the higher peaks are
known only to within about one copy (15\,Hz). They should be quoted as
approximate until the timing is fixed at acquisition.

\section{The board keeps up; our loop does not}

The guide's verification test used the vendor script \code{test\_modbus.py}
with the same 84-set maximum read: 448\,517 samples in 300\,s, 0 overruns,
0 CRC errors. Each batch uses about 22\,ms of the 67\,ms cycle on the serial
line. The rest is spent in our loop (\code{collect\_accel\_for\_duration} in
\code{main.py}). This supports the previous report's recommendation to make
that loop faster.

\section{Dropped samples are not buffered}

The guide describes a 1024-set buffer on the board, and our capture clears
it before each run. At about 16.5 surplus samples per read, the buffer
should take about 62 reads ($\approx$4.2\,s) to fill. Until then, no samples
should be lost: each batch would continue exactly where the previous one
ended, only later.

We tested this on the 30 flap ringdowns, most of which start between 3.0
and 4.5\,s. For the 62.7\,Hz mode we measured the phase change from one batch
to the next:

\begin{center}
\begin{tabular}{lc}
  \toprule
  & Phase change per batch \\
  \midrule
  Expected if batches are continuous & $-3.07$\,rad \\
  Expected if 67\,ms passes between batches & $+1.4$\,rad \\
  Measured, 2--4\,s (40 batch pairs) & $+1.09$\,rad \\
  Measured, 4--6\,s (74 batch pairs) & $+1.09$\,rad \\
  \bottomrule
\end{tabular}
\end{center}

The batches are not continuous even at 2--4\,s, so the samples are lost
from the start of the capture rather than after the buffer fills. Our
code does not lose them: it reads $\min(\text{AVAIL}, 84)$ sets and writes
all of them. Possible causes on the board:

\begin{enumerate}
  \item \textbf{A read discards the samples it does not return} (most
        likely). When more than 84 sets are waiting, the board returns 84 and
        discards the rest. This explains full batches, losses from the start
        and no backlog. It also explains the vendor's clean test, where fewer
        than 84 sets are ever waiting. OVERRUN would stay at 0.
  \item \textbf{The buffer is much smaller than 1024 sets.} It would fill
        within a fraction of a second, before any ringdown starts, so we
        cannot test this from existing data. OVERRUN would climb.
  \item \textbf{Samples are lost between the sensors and the board's
        processor} while it is busy answering a read. OVERRUN would stay at 0,
        and COUNTER would grow more slowly than 1493 per second. This is
        least likely, because the vendor's test handled the same data rate
        without loss.
\end{enumerate}

The guide is also out of date in at least one respect: it lists $\pm$2\,g
as the only full-scale range, but these ringdowns were recorded at
$\pm$16\,g. The buffer may likewise differ from what the guide describes.

We also checked whether each batch lines up with the exact time of its
read, which would show as the phase change following the $\pm$4\,ms jitter
in read times. The test was inconclusive: a phase estimated from 84 samples
(3.5 cycles) is too noisy. Either way, a batch sits at a roughly fixed
offset from its read time, so the corrected method's assumption holds.

\section{Recommendations}

\begin{enumerate}
  \item \textbf{Log AVAIL, OVERRUN and COUNTER on every read.} Our loop
        already reads AVAIL and discards it. One 10\,s capture tells the
        three causes apart: (1)~AVAIL $\approx$100 every read and OVERRUN
        $=0$; (2)~AVAIL levels off at the buffer size and OVERRUN climbs;
        (3)~COUNTER grows more slowly than $\approx$1493 per second. The
        logs also record exactly where each future capture's gaps are.
  \item \textbf{Read before 84 sets are waiting.} If cause~1 is right, this
        stops the losses without making the loop faster: for example, read
        smaller batches more often, so that AVAIL stays below 84, as in the
        vendor's test.
  \item \textbf{Use the measured rate, not 1600\,Hz,} in
        \code{plot\_ringdown\_psd.py}, \code{extract\_run\_summary.py} and the
        previous report's \code{make\_figures.py}. For future captures,
        estimate it per capture from COUNTER and the elapsed time. The
        sensors' rates are not crystal-controlled and may differ slightly
        from board to board.
  \item \textbf{Ask InUnit} how a read behaves when more than 84 sets are
        waiting, and what the actual buffer size is in the current firmware.
\end{enumerate}

\paragraph{Reproducing.} Run from the repository root:
\code{odr\_sweep.py} (Tables~\ref{tab:copies} and~\ref{tab:peaks}) and
\code{batch\_continuity.py} (the phase table), both in
\code{reports/pcb\_manual\_notes}.

\end{document}
